\section{The thermal source: heating and cooling}
\label{sec:thermal}

\subsection{The heating channels}
\label{sec:thermal:heat}

The volumetric heating rate of a cell is a sum of physically distinct
deposits, and there is exactly one list of them
(\texttt{util\_ion\_eq.f90}): \texttt{heating\_of\_composition} fills every
column of \texttt{heat\_chan}, the total it returns is the sum of those
columns in that order, and the header of \texttt{output/Heating\_breakdown.txt}
is built from the same names, so a deposit cannot exist in the energy
equation and be missing from the breakdown.

\begin{center}
\begin{tabular}{@{}>{\raggedright\arraybackslash}p{0.075\textwidth}>{\raggedright\arraybackslash}p{0.40\textwidth}>{\raggedright\arraybackslash}p{0.44\textwidth}@{}}
\toprule
\# & name & what it is\\
\midrule
1--6 & \texttt{heat\_HI \dots\ heat\_metals} &
  photoionization of each absorber, $n_a h^{(1)}_a$ of
  \eqref{eq:rad:photoheat}--\eqref{eq:rad:heatcontraction}\\
7  & \texttt{heat\_Hpe[excitedH]} & photoelectric heating of H($n{=}2$) by
     the Balmer continuum, \eqref{eq:ion:balmer}\\
8  & \texttt{heat\_Hdx[Lya-deexc]} & collisional de-excitation of the
     Ly$\alpha$-pumped $n=2$ atoms (\texttt{incl\_deexc\_heat}, default on)\\
9  & \texttt{heat\_He\_recomb} & photoelectrons of the H\,{\sc i}, H$_2$ and
     metal ionizations driven by helium recombination radiation\\
10--12 & \texttt{heat\_He23S\_*} & He($2^3$S)$+$H Penning and associative
     branches, and He($2^3$S)$+$H$_2$ Penning\\
13--14 & \texttt{heat\_H2\_LW\_*} & kinetic energy of the Lyman-Werner
     dissociation fragments, and collisional de-excitation of the H$_2$
     fluorescence\\
15 & \texttt{heat\_mol\_chem} & net chemical heat of the collisional
     reactions of the H$_2$/He network\\
16--19 & \texttt{heat\_FUV\_photolysis}, \texttt{heat\_oxygen\_collisional},
     \texttt{heat\_CO\_Hep\_transfer}, \texttt{heat\_CO\_photodissoc} &
     the oxygen option's deposits\\
\bottomrule
\end{tabular}
\end{center}

Channels 1--6 are the photoheating integrals of
Sec.~\ref{sec:radiation:rates}, already carrying the heating efficiency
$f_{\rm heat}$ of \eqref{eq:rad:fheat}. Channel 15 exists because the return
path that makes ionization energy a non-channel in an atomic gas is different
in a molecular one (\texttt{molecular\_reaction\_heat.f90}). Follow one
closed cycle:
\[
{\rm H_2}+h\nu \to {\rm H_2^+}+e \quad(\text{the photon pays }
I({\rm H_2})=15.43\,{\rm eV}),
\qquad
{\rm H_2^+}+{\rm H_2}\to{\rm H_3^+}+{\rm H} \quad(+1.70\,{\rm eV}),
\]
\[
{\rm H_3^+}+e \to {\rm H_2}+{\rm H}\quad(+9.25\,{\rm eV},\ \text{dissociative:
the energy goes into the fragments, not a photon}).
\]
The cycle returns to H$_2$ having turned one H$_2$ into H$+$H and having left
$I({\rm H_2})-D_0({\rm H_2}) = 15.43-4.48 = 10.95$\,eV in the gas, with no
photon anywhere in it to carry that away. Measured on the converged
He/H $=0.0793$ hot Uranus, that is 81 percent of the total heating rate at
$1.02\,r_{\rm base}$. \texttt{Molecular reaction heat} is therefore
\emph{on} by default; it is unreachable in an atomic run, so no non-molecular
result moves.

\subsection{The cooling channels}
\label{sec:thermal:cool}

\texttt{eval\_cool} (\texttt{util\_ion\_eq.f90}) assembles the cooling rate
of every cell as
\begin{equation}
\Lambda = n_e\left(\Lambda_{\rm ff}+\Lambda_{\rm cx}
                 +\Lambda_{\rm rec}+\Lambda_{\rm ci}\right)
        + \Lambda_{\rm metal} + \Lambda_{\rm H_3^+}
        + \Lambda_{\rm H_2} + \Lambda_{\rm H_2O} + \Lambda_{\rm CO},
\label{eq:th:cooltotal}
\end{equation}
in erg\,cm$^{-3}$\,s$^{-1}$, with $n_e$ the same charge sum the ionization
balance uses (\texttt{calc\_ne}), molecular donors included, since inside a
deep molecular base H$_3^+$ is the dominant ion. The breakdown written to
\texttt{output/Cooling\_breakdown.txt} is read straight off the arrays
\eqref{eq:th:cooltotal} is built from, so the columns sum to $\Lambda$
exactly.

\paragraph{Recombination.} $\Lambda_{\rm rec} =
\lambda_{\rm HII}n_{\rm HII}+\lambda_{\rm HeII}n_{\rm HeII}
+\lambda_{\rm HeIII}n_{\rm HeIII}$ with, for hydrogen \citep{Hui1997},
\begin{equation}
\lambda_{\rm HII}(T) = 3.435\times10^{-30}\,T\,
   \frac{x^{1.970}}{\left[1+(x/2.250)^{0.376}\right]^{3.720}},
\qquad x = \frac{2\times157807}{T}.
\label{eq:th:reccool}
\end{equation}
The helium terms are $kT$ times the \emph{same} recombination coefficient
the ionization balance runs on, so electrons are not removed at one rate and
the gas charged at another.

\paragraph{Collisional ionization.} $\Lambda_{\rm ci}$ charges each
ionization its measured potential in erg,
\begin{equation}
\Lambda_{\rm ci} = E_{\rm th}^{\rm HI}\beta_{\rm HI}n_{\rm HI}
   + E_{\rm th}^{\rm HeI}\beta_{\rm HeI}n_{1^1{\rm S}}
   + E_{\rm th}^{\rm HeII}\beta_{\rm HeII}n_{\rm HeII}
   + E_{\rm th}^{2^3{\rm S}}\beta_{2^3{\rm S}}n_{2^3{\rm S}},
\label{eq:th:cicool}
\end{equation}
the He\,{\sc i} term acting on the ground \emph{singlet}
$n(1^1{\rm S}) = n({\rm He\,I}) - n(2^3{\rm S})$, because the metastable
sits 19.8\,eV up and carries its own 4.8\,eV threshold in the last term.
The same subtraction is made in the photoheating and in the collisional
excitation below, so the metastable is never charged the ground-state
coefficients as well.

\paragraph{Collisional excitation.} $\Lambda_{\rm cx}$ is the sum of three
fits,
\begin{equation}
\lambda_{\rm HI} = \frac{7.5\times10^{-19}}{1+\sqrt{T/10^{5}}}e^{-118348/T},
\qquad
\lambda_{\rm HeI} = 1.1\times10^{-19}T^{0.082}e^{-2.3\times10^{5}/T},
\qquad
\lambda_{\rm HeII} = \frac{5.54\times10^{-17}T^{-0.397}}{1+\sqrt{T/10^{5}}}
                     e^{-473638/T},
\label{eq:th:coexcool}
\end{equation}
contracted with $n_{\rm HI}$, $n(1^1{\rm S})$ and $n_{\rm HeII}$. The first
of these is the Ly$\alpha$-dominated hydrogen line cooling, and it is the
optically thin fit as written: \emph{no escape probability multiplies the
Ly$\alpha$ cooling channel.} $\lambda_{\rm HI}$ is
\texttt{lambda\_coex\_HI} of \texttt{Cool\_coeff.f90} and
\texttt{eval\_cool} contracts it with $n_{\rm HI}$ and $n_e$ with no such
factor anywhere. The two escape probabilities the code does compute enter
elsewhere: $\beta$ of \texttt{lya\_rt.f90} \eqref{eq:ion:neufeld} enters the
H($n{=}2$) populations and the pumping field, and $\beta_{ul}$ of
\texttt{fine\_structure\_line\_transfer} enters the metal ground-term
statistical equilibrium \eqref{eq:th:fsse}; neither is read by any cooling
coefficient. The asymmetry is deliberate: multiplying $\lambda_{\rm HI}$ by
an escape probability would leave the trapped fraction with no owner while
keeping the de-excitation heating channel 8 as an independent source, that is
it would double count. A ledger carrying both is the stated target and is not
built.

With the metastable tracked, two further terms are added to
$\Lambda_{\rm cx}$, both scaling with the explicitly computed
$n(2^3{\rm S})$: the $2^3$S$\to2^3$P excitation followed by 10830\,\AA{}
decay and escape, and the $2^3$S$\to2^1$S (0.80\,eV) and $2^3$S$\to2^1$P
(1.40\,eV) collisional conversions, each removing its threshold energy from
the electron gas. The ground-to-triplet excitation $q_{13}$ is deliberately
excluded here, being already inside $\lambda_{\rm HeI}$.

\paragraph{Free-free.} $\Lambda_{\rm ff} = 1.426\times10^{-27}\sqrt{T}
\sum_{\rm ions} Z_{\rm ion}^{2}\,\bar g(T,Z_{\rm ion})\,n_{\rm ion}$, with
$\bar g$ the thermally averaged Gaunt factor of \citet{vanHoof2014},
interpolated linearly in $x = \log_{10}(Z^{2}\,157807/T)$
and clamped to the tabulated range. Only $Z_{\rm ion}=1$ and 2 occur.
The molecular ions H$_2^+$, H$_3^+$ and HeH$^+$ do donate to $n_e$ but are
deliberately left out of this charge sum: they exist only in the
$T\sim10^{3}$\,K molecular base, where this hot-plasma expression is an
extrapolation whose emission comes out at frequencies the atmosphere is not
thin to, and where H$_3^+$ and metal lines dominate by far.

\paragraph{Metal lines.} $\Lambda_{\rm metal} = n_e\sum_i n_{{\rm m},i}
c_i(T)$ with $c_i$ a closed-form analytic fit to the CHIANTI v11.0.2
optically thin line-cooling curve of that ion. The C/N/O fits are the
default (\texttt{cno\_chianti = .true.}; \texttt{cno\_cool 0} in
\texttt{metals.inp} restores the legacy set ported from AIOLOS), and
Mg\,{\sc i/ii}, Ca\,{\sc ii}, Na\,{\sc i} and Fe\,{\sc i/ii} follow the same
route. Two entries are not in CHIANTI v11 and are built otherwise: the
Na\,{\sc i} D doublet with a $\bar g$ factor 0.2 calibrated against the
CHIANTI Mg\,{\sc ii} and Ca\,{\sc ii} values, and the Fe\,{\sc i} permitted
multiplets.

Two density-dependent replacements are made on top of those coronal curves,
because a coronal fit assumes $\Lambda\propto n_e$ and no level saturation
while the base of an irradiated atmosphere runs at $n_e\sim10^{9}$,
$n_{\rm HI}\sim10^{14}$\,cm$^{-3}$, decades above the critical densities of
the lines involved.

\begin{itemize}
\item \emph{Fe\,{\sc ii}} takes a two-dimensional statistical-equilibrium
  table $\Lambda_{\rm eff}(T,n_e) = \left(\sum_u n_u A_{ul}\Delta E_{ul}\right)/n_e$,
  whose $1/n_e$ cancels the $n_e$ of the assembly and so leaves the correct
  saturated emission; at low $n_e$ it reduces to the coronal rate. Without
  it the forbidden $a\,^6$D lines overestimate the base cooling by about
  four decades.
\item \emph{Ground-term fine structure} of C\,{\sc i} ([C\,{\sc i}] 609.1,
  370.4\,$\mu$m), C\,{\sc ii} (157.7\,$\mu$m), N\,{\sc ii} (205.3,
  121.8\,$\mu$m) and O\,{\sc i} (63.2, 145.5, 44.1\,$\mu$m) is replaced by
  the exact statistical equilibrium of the ground term,
  \begin{equation}
  \sum_j f_j R_{ji} = f_i\sum_j R_{ij},
  \quad \sum_i f_i = 1,
  \quad R_{ul} = C_{ul}+\beta_{ul}A_{ul},
  \quad R_{lu} = C_{ul}\frac{g_u}{g_l}e^{-E_{ul}/kT},
  \label{eq:th:fsse}
  \end{equation}
  with $C_{ul} = n_e k_{e,ul}(T)+n_{\rm HI}k_{{\rm H},ul}(T)$ and
  $k_{e,ul} = 8.629\times10^{-6}\Upsilon/(g_u\sqrt T)$, the effective
  collision strengths $\Upsilon$ being quartic fits in
  $x = \log_{10}(T/10^{4})$ to the CHIANTI \texttt{.scups} values over
  $10^{3}$--$10^{5}$\,K (largest error 4.8 percent) and the H-atom rates
  taken from \citet{YanBabb2023} for C\,{\sc i} and N\,{\sc ii},
  \citet{Abrahamsson2007} for O\,{\sc i} and \citet{Barinovs2005} for
  C\,{\sc ii}. The
  emitted power per ion is
  $W_{\rm FS} = \sum_{u>l} f_u\beta_{ul}A_{ul}k_{\rm B}E_{ul}$ and the
  coefficient returned is $\Lambda_{\rm eff} = W_{\rm FS}/n_e+\Lambda_{\rm rem}$,
  $\Lambda_{\rm rem}$ being a refit of the coronal curve keeping only the
  channels that leave the ground term.
\end{itemize}

The line trapping $\beta_{ul}$ of these fine-structure lines is the
line-center escape probability formed from the columns above and below the
cell (\texttt{fine\_structure\_line\_transfer}), and it enters \emph{inside}
the statistical equilibrium as $A_{ul}\to\beta_{ul}A_{ul}$ rather than as a
factor on the result: in the subcritical limit the cooling is set by the
collisional excitation rate and is independent of $\beta$, and in the
saturated limit the escaping flux is proportional to it. $\beta = 1$
reproduces the optically thin result exactly, and that is the value every
other metal ion keeps. With \texttt{Base IR field: True} (default off) the
same lines also \emph{absorb} the thermal infrared of the lower atmosphere,
a blackbody at $T_0$ covering the sky fraction the base subtends, through
the photon occupation number $\bar n_{ul}$; this gives the layer a radiative
equilibrium floor instead of letting it cool into vacuum.

\paragraph{Molecular infrared.} $\Lambda_{\rm H_3^+}$ is the
\citet{Miller2013} emission of one H$_3^+$ molecule in LTE,
\begin{equation}
\Lambda_{\rm H_3^+} = 4\pi\,n_{\rm H_3^+}\,E(T)\,s(T,n_{\rm H_2}),
\qquad
\ln E(T) = \sum_n C_n T^{n}\ \ \mathrm{[W\,molecule^{-1}\,sr^{-1}]},
\label{eq:th:h3p}
\end{equation}
their Table 5 (the ``using $z(t)$'' column) with the non-LTE departure
factor $s$ of their Table 6, whose only collider is H$_2$ through the
proton-hopping reaction with the \citet{Oka2004} coefficient
$2\times10^{-15}$\,m$^{3}$\,s$^{-1}$. Table 6 covers
$n({\rm H_2}) = 10^{6}$--$10^{14}$\,cm$^{-3}$; below that the code adopts
the linear low-collider limit $s(T,n) = s(T,10^{6})\,n/10^{6}$ and records
the cell as evaluated outside the tabulated range. This is the dominant
coolant of a molecular base: the atomic channels above are all exponentially
suppressed at $T\sim10^{3}$\,K, so leaving it out leaves that gas with no
radiative loss at all.

With \texttt{Molecular IR bands: True} (default off) three further net
exchanges are added (\texttt{molecular\_infrared\_cooling.f90}): the H$_2$
quadrupole and magnetic-dipole line spectrum, and the H$_2$O and CO
vibration-rotation bands. Each is a \emph{net} rate,
\begin{equation}
\Lambda_X^{\rm net} = n_X\left[E_X(T)-W_{\rm dil}\,E_X^{\rm abs}(T;T_{\rm rad})\right],
\qquad
E_X(T) = 4\pi\sum_b\sigma_b(T)\int_b B_\nu(T)\,{\rm d}\nu,
\label{eq:th:molir}
\end{equation}
so each vanishes at its own radiative equilibrium temperature instead of
running the layer down to nothing; a column is negative wherever that channel
heats. The dilution $W_{\rm dil}$ is the same $0.5\,f_{\rm sky}$ the H$_3^+$ closure
uses and is zero when \texttt{Base IR field} is off, which reduces these
channels to pure emitters. Without them a converged molecular layer radiates
itself down to 190--400\,K against an equilibrium temperature of
1100--1400\,K.

\subsection{Advancing the temperature}
\label{sec:thermal:tsolve}

\paragraph{In the marching loop.} The temperature of a cell is advanced by
\texttt{solve\_energy\_semi\_implicit}
(\texttt{energy\_semi\_implicit.f90}), the inner step of the coupled
temperature-composition source step: the marching loop alternates it with
the composition sweep until the pair stops moving, so the composition is
fixed for one pass and not for the whole step. With the composition and the
photoheating frozen over the step,
\begin{equation}
\rho e(T)-\rho e_{\rm old} = \Delta t\left[\mathcal{H}-\Lambda(T)\right],
\qquad\text{i.e.}\qquad
R(T) = \varepsilon(T)-\varepsilon_{\rm old}
       -a\left[\mathcal{H}-\Lambda(T)\right] = 0,
\quad a = \frac{\Delta t}{n_{\rm tot}+n_e},
\label{eq:th:semiimplicit}
\end{equation}
with $\varepsilon = \rho e/(n_{\rm tot}+n_e)$ the internal energy per
particle, which is what \texttt{internal\_energy\_per\_particle} returns.
The caloric internal energy $\rho e$ is that of \eqref{eq:caloric}. The root is
found by a safeguarded Newton-bisection hybrid on a physical bracket whose
ends are the floor of the equation-of-state validity and a ceiling above the
heating-only update: the search starts from the explicit-Euler point,
brackets the root as it goes, and once both ends are known takes a Newton
step whenever that step stays strictly inside the bracket and the bracket
keeps halving, bisecting otherwise. The Newton derivative is exact in the
equation-of-state term (\texttt{heat\_capacity\_per\_particle}) and a secant
estimate in ${\rm d}\Lambda/{\rm d}T$, which affects only the path, the
safeguard owning the convergence.

What is tested before the update is returned is the residual at the
\emph{returned} temperature against $|R|/s$ with
$s = |\rho e_{\rm old}|+\Delta t(|\mathcal{H}|+|\Lambda|)$, and admissibility; the
cooling written back is $\Lambda(T_{\rm returned})$, evaluated there and not
at a previous iterate. There is no clamp: a cell that does not meet the
tolerance, that has no bracket, or that reaches the equation-of-state floor
is a named failure, nothing is assembled from it, and the step controller
retakes the step at half $\Delta t$ as in Sec.~\ref{sec:numerics:rk}. The floor is
not a reservoir: reaching it means the balance asks for a temperature below
the range in which the equation of state and the rate fits are defined, and a
floor representing real external heating would have to be a specified
reservoir with its own budget. \texttt{use\_semi\_implicit\_energy} is on by
default.

\paragraph{In the advection post-process.} \texttt{T\_equation.f90} is
\emph{not} the marching loop's temperature solve. Its module
\texttt{equation\_T} is used by \texttt{post\_process\_adv.f90}
(Sec.~\ref{sec:postprocessing:adv}) and by test drivers, and by nothing else;
what it solves is the \emph{steady} internal-energy balance of the profile,
upwind differenced,
\begin{equation}
\nabla\!\cdot\!(u\mathbf{v}) + p\,\nabla\!\cdot\!\mathbf{v}
 = \rho v\frac{{\rm d}e}{{\rm d}r}
 - p\,v\frac{{\rm d}\ln\rho}{{\rm d}r}
 + h\,\nabla\!\cdot\!(\rho\mathbf{v})
 = \mathcal{H}-\Lambda,
\label{eq:th:teq}
\end{equation}
with $e$ the internal energy per unit mass and $h = e+p/\rho$ the enthalpy
per unit mass. The upstream term carries the \emph{upstream} cell's specific internal
energy formed at the upstream composition and temperature, not this cell's
caloric state evaluated at $T_{\rm up}$: the rovibrational heat capacity
belongs to the gas that holds the molecules, so across a dissociation front
two cells at the same temperature store different energy. The root is
bracketed rather than guessed (\texttt{solve\_T\_brent}, on by default
through \texttt{use\_brent\_tsolve}): 80 logarithmically spaced probes from
$\max(0.05T_{\rm guess},1\,{\rm K})$ to $4T_{\rm guess}$ locate the
\emph{first} sign change, i.e.\ the lowest root, and Brent's method closes
the bracket to $2\epsilon|b|+{\rm tol}/2$ with ${\rm tol}=10^{-10}$ in
$T/T_0$, which is what keeps the solve off the hot spurious root;
\texttt{hybrd1} on the same residual is the fallback. The metal densities and
the fine-structure escape probabilities are held fixed while the root finder
varies $T$, the latter because the optical depth is a column over the whole
atmosphere and not a function of this cell's trial temperature.

\paragraph{The heating efficiency.} The diagnostic $\eta_{\rm heat}$ of
\eqref{eq:rad:eta}, written as column \texttt{eta} by
\texttt{write\_output.f90}, is a property of the composition the sweep
returned and is not a free parameter anywhere in the code; the closed-form
parameterizations of the same quantity that the literature uses are
reproduced offline by \texttt{eta\_approx.py} for comparison only.
